\section{Introduction}\label{sec:intro}

\revised{This section describes Accolade's pricing problem, reviews related work, defines the Phase~1
task, and outlines our approach.}

\subsection{Motivation and business problem}

\revised{Accolade, a property management platform, is adding a revenue management product to its
offerings.} The goal is to help property owners price their units
correctly: an apartment priced too low leaves revenue on the table, while one priced too high sits
vacant, and every vacant day is rent the owner does not recover.

\revised{The right price depends on what comparable units are asking now, and finding those comparables is
a challenge in itself. Nearby units are not always comparable, because unit characteristics,
amenities, building quality, concessions, and fees also affect price.}

\revised{This report distinguishes three rents.} The asking rent is
the price posted on a listing. The effective rent is the asking rent less any concessions,
such as a free month, spread over the lease:
\[
  \text{effective rent} \;=\; \text{asking rent} \;-\; \text{amortized monthly concession}.
\]
The market rent is what a unit is actually worth given the market around it. Accolade's
product aims to estimate the third from listings that only show the first.

\revised{Accolade's Draft Project Plan builds the product in two phases.} The first, a
market rent estimate, determines what comparable properties are currently asking for a given
unit type. The second, a dynamic price recommendation, determines the asking price a property
should post today based on its current price, demand relative to its occupancy target, comparable
prices, price sensitivity, and recent leases. Phase~2 builds on Phase~1, since the recommendation
uses the comparable set and the market rent estimate as its anchors.

\subsection{Existing work}

Rent estimation has long been approached through hedonic pricing, which treats the price of a
property as the sum of the values of its characteristics~\autocite{rosen1974hedonic}. To ground our
approach in more recent work, each member of the team reviewed a paper covering a modeling
technique, the effects of algorithm-based pricing, or the measurement of rental markets from
listing data.

\revised{A study of residential property valuation in Germany treated comparable-property selection as the
baseline that learned models must beat: it compared an expert function for selecting comparable
properties with hedonic models (OLS and GAM) and XGBoost, and found that XGBoost predicted
best~\autocite{stang2023human}.} Other studies show the value of richer inputs: a
Melbourne study combined text, image, and geospatial embeddings~\autocite{hasan2024multimodal}, a São
Paulo study paired location-based features with natural language processing of listing
descriptions~\autocite{maia2024machine}, and a reproducible study of the major Spanish markets used
machine learning to identify spatial market segments~\autocite{reyblanco2024using}. Finally,
Zillow's home value index methodology shows how \revised{to measure values consistently over
time}~\autocite{olsen2023zhvi}, while a \revised{study of algorithmic pricing in U.S. multifamily rentals
estimated} its effects with two-way fixed effects (TWFE) models~\autocite{calderwang2026algorithmic}.

These ideas map directly onto Accolade's data: automated comparable selection onto the platform's
suggested comparables, image embeddings onto its listing photos and quality scores, and time-based
validation onto its daily price series. \revised{They also guide our choice of models}, from hedonic
regression to Random Forest, XGBoost, and stacked ensembles.

\subsection{Problem statement}

The technical goal of Phase~1 is to estimate the asking rent of a unit type, defined by its
bedrooms, bathrooms, and a square-footage band, as a point estimate together with a range. \revised{Every
estimate must meet two requirements. First, it must include its range and the number of comparables
or observations behind it.} Second, every price must
be explainable in terms of which comparables were used, which adjustments were applied, and which
limits bound the result. \revised{Net effective rent is a second, derived output,
which we will produce once concessions can be reliably extracted from listing text.}

This report covers the Phase~1 modeling completed so far\revised{, our first models on real Dwellsy
listings,} and the first component of Phase~2, the extraction of concession terms from advertised
specials.

\subsection{Overview}

\Cref{fig:pipeline} summarizes our pipeline. We merged the source tables of the synthetic dataset
into a single master table with one row per listing, screened its variables for information that is
not available when a unit is listed, scored the remaining variables, and trained and compared models
on buildings held out from training. The generator's truth layer is used only to evaluate the
predictions, never as an input. \revised{We have begun applying the same approach to real Dwellsy
listings from Chicago.} \Cref{sec:methods} describes the data, preprocessing, exploratory
analysis, and methods; \cref{sec:results} presents our preliminary results; and
\cref{sec:discussion} discusses our conclusions and next steps.

\begin{figure}[htbp]
  \centering
  % Pipeline: source layers to evaluation. Numbers come from generated/numbers.tex.
\begin{tikzpicture}[
  font=\scriptsize, >={Stealth[length=1.7mm, width=1.4mm]},
  box/.style={draw=rulecolor, line width=0.6pt, fill=white, align=flush left, anchor=north west,
              inner xsep=5pt, inner ysep=4pt},
  layer/.style={box, fill=wash, text width=2.45cm},
  stage/.style={box, text width=2.5cm, minimum height=2.05cm},
  flow/.style={->, draw=ink2, line width=0.7pt},
  h/.style={font=\scriptsize\bfseries\color{ink}},
]
\node[layer] (market) at (0, 0) {\textbf{Market}\\\textcolor{ink2}{listings, buildings,\\specials}};
\node[layer, below=1.8mm of market] (operator) {\textbf{Operator}\\\textcolor{ink2}{leases, leads\\(Phase 2)}};
\node[layer, below=1.8mm of operator, draw=brick, fill=white] (truth) {\textbf{\textcolor{brick}{Truth}}\\\textcolor{ink2}{true rents,\\answer keys}};
\node[anchor=south west, inner sep=0, text=ink2] at (0, 0.12) {\nSourceTables{} tables\enspace·\enspace\nSourceRowsM\,M rows};

\node[stage] (master) at (3.4, 0) {\textbf{Master table}\\[2pt]\textcolor{ink2}{one row per listing\\\nMasterRows{} rows\\\nMasterCols{} columns}};
\node[stage] (select) at (6.7, 0) {\textbf{Feature selection}\\[2pt]\textcolor{ink2}{roles screen leaks,\\then scoring\\\nMasterCols{} \(\to\) \nFeatures{} inputs}};
\node[stage] (models) at (10.0, 0) {\textbf{Models}\\[2pt]\textcolor{ink2}{baseline, linear,\\decision tree,\\gradient boosting}};
\node[stage, draw=accent, line width=0.8pt] (eval) at (13.3, 0) {\textbf{\textcolor{accent}{Evaluation}}\\[2pt]\textcolor{ink2}{\nTestBuildings{} unseen\\buildings, scored\\against asking\\and true rent}};

\draw[flow] (market.east) -- (master.west |- market.east);
\draw[flow] (master.east) -- (select.west);
\draw[flow] (select.east) -- (models.west);
\draw[flow] (models.east) -- (eval.west);
\draw[flow, draw=brick] (truth.east) -- (eval.south |- truth.east) -- (eval.south);
\node[fill=white, inner xsep=3pt, text=brick] at ($(truth.east)!0.5!(eval.south |- truth.east)$)
  {scores predictions, never an input};
\end{tikzpicture}

  \caption{From source tables to evaluation. The generator's truth layer feeds only the evaluation:
  it scores the predictions and is never a model input.}
  \label{fig:pipeline}
\end{figure}
